\section{Related Work}

\subsection{Physarum models of shortest path}

The slime mould \emph{Physarum polycephalum} connects food sources through a
tube network that thickens where flux is high and atrophies where it is
low, and it famously resolves a maze into the shortest connecting
path~\cite{nakagaki2000maze}. Tero et al.\ formalised this as a system of
adapting conductivities driven by a Poiseuille-flow pressure
solve~\cite{tero2007mathematical}; the same group later used it to produce
networks comparable to the Tokyo rail system~\cite{tero2010rules}. Bonifaci
et al.\ proved that the continuous dynamics converge to the shortest path,
for \emph{undirected} graphs~\cite{bonifaci2012physarum}. We implement the
canonical Tero model, with the simpler Bonifaci adaptation rule available as
a configuration switch. Awad et al.\ survey the wider family of
Physarum-inspired applications~\cite{awad2022survey}.

The directedness restriction is explicit in the literature rather than
incidental. Zhang et al.\ note that the model is designed only for
undirected networks and propose a directed formulation based on analog
circuit analysis~\cite{zhang2014directed}. We do \emph{not} implement that
variant; the behaviour we report for the canonical model on a directed
network should therefore be read as a measurement of an acknowledged
limitation, not as the discovery of a defect.

\subsection{The prior three-way comparison}

The combination of Physarum, ACO and Dijkstra in one study is not new.
Zhang et al.\ benchmark exactly these three~\cite{zhang2014improved} and
report that their improved Physarum algorithm and Dijkstra both attain
optimal paths ``with one hundred percent'', while ACO's accuracy degrades
with network size. Their evaluation uses Erd\H{o}s--R\'enyi random graphs of
15 to 2{,}000 nodes with static integer edge lengths; there is no congestion
model, no disruption and no real road network.

Two aspects of that paper shape ours. First, \emph{it needs no extraction
machinery}: the path is read directly off the converged flux field, because
on its test graphs the flux separates cleanly, going to one along the path
and zero elsewhere. It therefore has nothing to pool, and we cannot claim
that Physarum evaluations in general conflate extraction provenance. The
honest contrast is that on random graphs extraction is a formality, whereas
on a real directed road network it becomes the step that decides the reported
result. Second, their Fig.~9 places ACO at zero optimal-path accuracy from
1{,}000 nodes upward. Our ACO, on a larger and considerably harder network,
attains 43--44\% exact match and 73.3--81.3\% within 10\% of optimal; we
note this not to claim superiority, since the graph families differ, but to
establish that the ACO baseline used here is a strong one by the standard of
the prior art.

\subsection{Extraction and termination}

Closest to our extraction finding is the transition-point analysis of Wang et
al.~\cite{wang2021transition}, which defines a ``dominant path'' procedure
--- repeatedly step to the maximum-flux edge, delete the visited node, and
continue --- that is procedurally identical to the maximum-flux walk we
implement. They state that when the solver converges, the dominant path is
exactly the shortest path of the graph. On our network it is not: at the
conductivity convergence point the walk yields a directed-legal path on
43.0\% of solves and matches the optimum on 23.3\%.

The two results are compatible, and the reason is the test substrate. Their
motivating examples are five-node \emph{complete} graphs, and their main
datasets are complete graphs of 50--5{,}000 nodes and small-world graphs of
mean degree 6 to 600, with the graph assumed undirected throughout. On a
complete graph the maximum-flux walk cannot dead-end, because every
unvisited node remains adjacent; the failure mode we measure is structurally
unreachable there. Our network has mean out-degree 2.46 and pervasive
one-way streets. We also adopt their methodological point --- that apparent
speed-ups of accelerated Physarum variants can stem from an inappropriate
termination criterion --- and test their termination rule on a substrate they
did not evaluate.

\subsection{Ant colony optimization on traffic networks}

MMAS~\cite{stutzle2000maxmin} bounds pheromone to $[\tau_{\min},
\tau_{\max}]$, which is what makes repeated re-adaptation possible; plain Ant
System stagnates. Among traffic-network applications, Zhang et al.\ propose
orientation guidance, weighting edges by their angular alignment with the
destination~\cite{zhang2018novel}. That is related to, but not the same as,
the goal-directed ablation we evaluate, which folds an A*-style great-circle
time bound into the visibility term. Their evaluation uses a single 528-node
network with static arc travel times and compares only against conventional
ACO, so our result --- that goal direction is neutral under free flow and
harmful under BPR-heavy congestion --- extends rather than contradicts it.

More recently, Tuan et al.\ report an improved ACO for dynamic shortest-path
problems that outperforms conventional ACO, Dijkstra and a genetic
algorithm~\cite{tuan2026improved}. Although road closures appear in its
motivation, the evaluation uses the METR-LA loop-detector
dataset~\cite{li2018diffusion} under fluctuating traffic, and does not
measure recovery following a closure. It is nevertheless the closest work to
ours in one respect that counts against us: its congestion is measured
rather than simulated. We return to this in Section~\ref{sec:limitations}.

\subsection{Position of this work}

We verified the above with a systematic pass over the open citation
databases, including complete forward-citation sweeps of both
\cite{zhang2014improved} (28 citing works) and \cite{watanabe2011traffic}
(69 citing works). That search returns no study combining all three
algorithm families on a real road network under simulator-generated
congestion with closure recovery as a measured outcome; in particular, no
work pairing a Physarum model with SUMO was found at all. Related but
distinct are Physarum/ACO \emph{hybrids}, which fuse the two mechanisms into
one algorithm rather than comparing them. The contribution claimed here is
therefore the substrate and the measurement --- a real, directed road
network, simulator congestion, and recovery --- and not the combination of
the three algorithms, which \cite{zhang2014improved} already reported.
